\section{Results}
\label{results}
\subsection{Experimental Setup}



To demonstrate general-purpose capability as a vision encoder, we evaluate our models on a suite of 9 standard tasks, covering two broad types: image-text and image-only.
Datasets and metrics used in this work are introduced in the following paragraphs.
Evaluations follow the detailed protocol of \citep{tips_paper}, to ensure fair comparisons.
As our objective is to assess the applicability of our model off-the-shelf for downstream applications, the pretrained image-text representations are kept frozen for all evaluations.

\noindent\textbf{Image-text tasks} include 2 types, focused on dense and global image-text alignment.
We consider dense image-text alignment by evaluating on zero-shot segmentation, reporting mean Intersection over Union (mIoU).
Unless otherwise mentioned, we report \tipsvtwo results on zero-shot segmentation with the simple protocol that maximizes cosine similarities between patch features, up-scaled to image resolution, and the text features of each class name.
Note that DINOv2 (with text~\cite{jose2024dinov2}) and SILC~\cite{naeem2024silc} instead report results using the expensive sliding window protocol from TCL~\citep{cha2023learning}, which generally boosts results.
Both protocols use the value embeddings of the final transformer layer as patch embeddings, following ~\citep{zhou2022extract}.
We report zero-shot segmentation on ADE20K~\citep{zhou2017ade20k} with 150 classes (ADE150), PASCAL context \citep{everingham2011pascal}, and PASCAL VOC~\citep{everingham2010pascal}.
For PASCAL context, we report both with (PC60) and without (PC59) the background label.
The global-focused tasks include image-to-text and text-to-image retrieval on Flickr30K~\citep{young2014from}, DOCCI~\citep{onoe2024docci}, and COCO~\citep{chen2015cococaptions} reporting recall@1; and zero-shot classification on ImageNet-1K reporting top-1 accuracy by retrieving nearest class text embeddings~\citep{radford2021clip}. 


\noindent\textbf{Image-only tasks} include semantic segmentation with linear probe on PASCAL VOC~\citep{everingham2010pascal} and ADE20k~\citep{zhou2017ade20k} reporting mIOU; monocular depth estimation on the scene-centric NYUv2~\citep{silberman2012indoor} with linear probe and the object-centric NAVI~\citep{jampani2023navi} with DPT~\cite{ranftl2021vision}, both reporting RMSE; surface normal estimation on NYUv2 and NAVI, using DPT, reporting angular RMSE; image classification on ImageNet-1K~\citep{russakovsky2015imagenet} reporting top-1 accuracy (KNN and linear probe); and fine-grained/instance-level retrieval on the Universal Embeddings Dataset (UnED)~\citep{ypsilantis2023uned} (comprising data from $8$ different domains: Food2k \citep{min2023large}, CARS196 \citep{krause20133d}, SOP \citep{song2016deep}, InShop \citep{liu2016deepfashion}, iNat \citep{van2018inaturalist}, Met\citep{ypsilantis2021met}, GLDv2 \cite{weyand2020google}, Rp2k\cite{peng2020rp2k}), reporting average recall@1.

\noindent\textbf{Training data.}
We train our models using the web-crawled WebLI dataset~\citep{chen23pali}, reusing the specific filtered subset created by~\cite{tips_paper}, containing $116$M images.
We also reuse synthetic captions from PaliGemma~\citep{beyer2024paligemma}, again following~\cite{tips_paper}.
Additionally, as described in \cref{multi_caption}, we leverage synthetic captions from Gemini 1.5 Flash \cite{geminiteam2024gemini}.

\noindent\textbf{Implementation details.}
Following \citep{tips_paper}, pretraining combines contrastive image-text and self-supervised losses, integrating the key upgrades described in \cref{sec:method}: iBOT++, head-only EMA and multi-granularity text captions.
Pretraining consists of a low-resolution stage for $90k$ steps at batch size $8192$ and then a high-resolution adaptation stage for $9k$ steps at batch size $4096$.
At low resolution, we train with $1$ global crop at resolution $224$ and $6$ local crops at resolution $98$.
At high resolution, the global crop has resolution $448$ and the local crops have resolution $140$.
We use only random resizes, crops and flips as augmentations.
We pre-train at ViT-g scale with 512 TPUv5 chips for 2 days.
We also distill to a family of -SO, -L, -B models, with all results presented in the appendix.
Distillation is as described in \cref{sec:preliminaries}, and then student models also undergo a high-resolution adaptation stage. 




\subsection{Ablations}
We present ablations in \cref{tab:ablation_studies} where we add each of the primary contributions successively, evaluated on a representative subset of 7 of our evaluation tasks.
We begin with TIPS~\citep{tips_paper} as a baseline, reproduced with two minor modifications: we halve the batch size and increase position embeddings to $32\times32$ from $16\times16$.
Our preliminary experiments showed halving the batch size saved training resources without affecting results significantly, and higher-resolution position embeddings improved all results without incurring much cost. 
Ablations are carried out for 100k steps at original resolution, without the high-res adaptation stage of the full pretraining process.
In the second row, iBOT is replaced by iBOT++, which results in across-the-board improvements for most evals, including a dramatic $+14.1$ mIoU improvement to zero-shot segmentation.
After that, we replace PaliGemma synthetic captions with multi-granularity captions sampled from PaliGemma and Gemini captions, which improves both global and dense types of image-text tasks: I$\rightarrow$T, T$\rightarrow$I, and zero-shot segmentation.
Finally, we add head-only EMA, a change that is primarily targeted at enabling significant training resource savings, but still results in comparable performance, even further improving some tasks such as  zero-shot segmentation. 
The combination of these changes is \tipsvtwo, with all primary components resulting in increased or comparable performance across our evaluations, with particular significant improvement for zero-shot segmentation from iBOT++. 


\begin{table}
\caption{\small\textbf{Dense image-text evaluations}, where \tipsvtwo outperforms others in all cases, even though SILC and DINOv2 use the more expensive TCL protocol~\cite{cha2023learning}. We highlight the \textbf{best} and \underline{second-best} number of each column.
}
\vspace{-2mm}
\scriptsize
\centering
{
\begin{tabular}{l c c c c}
\small \multirow{2}{*}{Method} & \multicolumn{4}{c}{Zero-shot Segmentation $\uparrow$}\\
& PC59 & PC60 & VOC21 & ADE150  \\
\midrule
SigLIP2 SO/14 & - & 19.6 & 26.8 & 15.6 \\
SILC B/16 & 31.6 & - & - & 19.3 \\
DINOv2 (dino.txt) L/14 & 30.9 & - & - & 20.6 \\
TIPS L/14 & \underline{33.5} & \underline{30.4} & \underline{30.5} & \underline{20.8} \\
\tipsvtwo L/14 (ours) & \textbf{37.1} & \textbf{33.9} & \textbf{44.4} & \textbf{24.7} \\

\end{tabular}
}

\label{tab:sota_dense_image_text}
\end{table}

\begin{table*}[ht]
\caption{\small\textbf{Global image-text evaluations}, where \tipsvtwo achieves best or second-best in $5$ out of $7$ cases. We highlight the \textbf{best} and \underline{second-best} number of each column.}
\vspace{-2mm}
\scriptsize
\addtolength{\tabcolsep}{0.5em}
\centering
{
\begin{tabular}{l c c c c c c c c}
\small \multirow{2}{*}{Model} & \multicolumn{3}{c}{I$\rightarrow$T Retrieval $\uparrow$} & \multicolumn{3}{c}{T$\rightarrow$I Retrieval $\uparrow$} & \multicolumn{1}{c}{0-Shot $\uparrow$} \\
& COCO & Flickr & DOCCI & COCO & Flickr & DOCCI & ImageNet \\
\midrule
CLIP L/14 & 56.3 & 85.2 & 44.4 & 36.5 & 65.2 & 40.4 & 75.5 \\
OpenCLIP G/14 & 67.3 & 92.9 & - & 51.4 & 79.5 & - & 80.1 \\
SigLIP2 g/16 & 72.8 & \underline{95.4} & - & 56.1 & \textbf{86.0} & - & \underline{85.0} \\
SILC G/16 & 73.2 & - & - & 54.7 & - & - & 83.7 \\
DINOv2 (dino.txt) L/14 & - & - & - & 45.4 & 77.1 & - & 81.4 \\
PE-core G/14 & \underline{75.4} & \textbf{96.2} & - & 58.1 & 85.7 & - & \textbf{85.4} \\
TIPS g/14 & 74.0 & 93.0 & \underline{57.2} & \underline{59.4} & 84.5 & \underline{58.8} & 79.9 \\
\tipsvtwo g/14 (ours) & \textbf{75.7} & 95.1 & \textbf{68.9} & \textbf{60.7} & \underline{85.9} & \textbf{72.8} & 80.7 \\
\end{tabular}
}
\vspace{-2mm}
\label{tab:sota_image_text}
\end{table*}
 
\begin{table*}[ht]
\caption{\small\textbf{Image-only evaluations}, where \tipsvtwo achieves the best or second-best performance in $7$ out of $9$ evaluations. We highlight the \textbf{best} and \underline{second-best} number of each column.
}
\vspace{-2mm}
\scriptsize
\addtolength{\tabcolsep}{-0.1em}
\centering
{
\begin{tabular}{l c c c c c c c c c }
\small \multirow{2}{*}{Method} & \multicolumn{2}{c}{Segmentation $\uparrow$} & \multicolumn{2}{c}{Depth $\downarrow$} &
\multicolumn{2}{c}{Normals $\downarrow$} & Fine-grained $\uparrow$ &
\multicolumn{2}{c}{ImageNet Class. $\uparrow$} \\
& PASCAL & ADE20k & NYUv2 & NAVI & NYUv2 & NAVI & retrieval (UnED) & KNN & lin \\
\midrule
CLIP L/14 & 74.5 & 39.0 & 0.553 & 0.073 & 24.3 & 25.5 & 57.4 & 79.8 & 84.3 \\
SigLIP2 SO/14 & 78.1 & 45.4 & 0.466 & 0.064 & 23.0 & 25.0 & - & - & - \\
DINOv2 g/14 & 83.1 & 49.5 & 0.372 & \bf{0.054} & \bf{20.7} & \bf{24.0} & 62.7 & 83.6 & \underline{87.3} \\
PE-core G/14 & - & 41.5 & - & - & - & - & - & \bf{86.8} & \bf{89.5} \\
PE-spatial G/14 & - & 49.3 & - & - & - & - & - & - & - \\
FRANCA G/14 & 81.3 & 42.4 & - & - & - & - & - & 83.0 & 85.9 \\
TIPS g/14 & \underline{83.6} & \underline{49.9} & \underline{0.353} & \underline{0.058} & 21.9 & 24.2 & \bf{68.2} & 83.3 & 86.2 \\
\tipsvtwo g/14 (ours) & \bf{85.1} & \bf{51.6} & \bf{0.334} & 0.059 & \underline{21.7} & \underline{24.1} & \underline{67.0} & \underline{83.7} & 86.8 \\
\end{tabular}
}
\vspace{-2mm}
\label{tab:sota_image_only}
\end{table*}

\begin{table*}[t!]
\caption{\small\textbf{Comparison between DINOv3 and \tipsvtwo}, on largest comparable model size ViT-L. \tipsvtwo achieves superior performance on 4 out of 6 metrics. We highlight the \textbf{best} performing model for each task.}
\vspace{-2mm}
\scriptsize
\addtolength{\tabcolsep}{0.5em}
\centering
{
\begin{tabular}{l c c c c c c}
\small \multirow{2}{*}{Model} & Segmentation $\uparrow$ & Depth $\downarrow$ & 0-Shot $\uparrow$ & I$\rightarrow$T Ret. $\uparrow$ & T$\rightarrow$I Ret. $\uparrow$ & Zero-shot Seg. $\uparrow$ \\
& ADE20k & NYUv2 & ImageNet & COCO & COCO & ADE150 \\
\midrule
DINOv3 L/16 & \textbf{54.9} & 0.352 & \textbf{82.3} & 63.7 & 45.6 & 24.7 \\
\tipsvtwo L/14 (ours) & 51.4 & \textbf{0.339} & 79.7 & \textbf{73.5} & \textbf{57.4} & \textbf{25.1} \\
\end{tabular}
}
\vspace{-2mm}
\label{tab:dinov3}
\end{table*}


\subsection{Comparisons with Other Vision Encoders}

We provide extensive benchmarking against state-of-the-art methods.
For fair comparisons, we strictly follow the protocol of \citep{tips_paper}, considering the largest instantiations in each model family, up to ViT size ``G", with $1.8$B parameters in the vision encoder.
Note that the largest model in the \tipsvtwo family is of size ``g", with $1.1$B parameters in the vision encoder.
The compared methods are the following.
Weakly-supervised methods: CLIP~\citep{radford2021clip}, PE~\citep{bolya2025perceptionencoderbestvisual}, and OpenCLIP~\citep{cherti2023reproducible}; self-supervised methods: DINOv2~\citep{oquab2024dinov2} (with registers~\cite{vitsneedregisters}), DINOv3~\citep{siméoni2025dinov3}, and FRANCA~\citep{venkataramanan2025franca}; combined approaches: TIPS~\citep{tips_paper}, SILC~\citep{naeem2024silc}, dino.txt~\citep{jose2024dinov2} (as the image-text capability of DINOv2), and SigLIP2~\citep{tschannen2025siglip2multilingualvisionlanguage}.
Qualitatively, \cref{fig:qualitative_dense} shows that \tipsvtwo patch embeddings are more spatially coherent than other vision-language pretraining methods.

\noindent\textbf{Dense image-text tasks.} In \cref{tab:sota_dense_image_text}, we report zero-shot segmentation results on ViT-L and the closest comparable sizes from competitor models.
\tipsvtwo achieves state-of-the-art results, indicating a significant improvement in dense patch-text alignment. 
Qualitative visualizations of the zero-shot segmentation results are presented in \cref{fig:zss_viz}.

\noindent\textbf{Global image-text tasks.} 
In \cref{tab:sota_image_text}, \tipsvtwo achieves the best or second-best results in 5 out of 7 evaluations. Our models are best in all COCO and long-text DOCCI tasks, outperforming even larger ViT-G models.
Note that our ViT-g models even outperform the recent PE~\citep{bolya2025perceptionencoderbestvisual} ViT-G models in 3 out of 5 evals, although the PE model has 56\% more parameters and processes 47$\times$ more image-text pairs.

\newcommand{\showpca}[2]{%
  \resizebox{#1\linewidth}{!}{%
	\setlength{\fboxsep}{0pt}
		\hfill
		\includegraphics[height=#1\linewidth]{images/#2.jpeg}
		\hfill
		\includegraphics[height=#1\linewidth]{images/tips/#2.png}
		\hfill
		\includegraphics[height=#1\linewidth]{images/siglip2/#2.png}
		\hfill
		\includegraphics[height=#1\linewidth]{images/supp/tipsv2/vit_g/#2.png}
}}

\newcommand{\pcawidth}{1.0}
\begin{figure}[tb]
	\centering
	\showpca{\pcawidth}{dadaocheng} \\[0.2mm]
    \showpca{\pcawidth}{cph} \\[0.2mm]
    \showpca{\pcawidth}{hike} \\[0.2mm]
    \showpca{\pcawidth}{angus} \\[0.1mm]

    \resizebox{\pcawidth\linewidth}{!}{%
    \scriptsize
        \begin{minipage}[t]{0.24\linewidth}
          \centering
          Image
        \end{minipage}%
        \hfill
        \begin{minipage}[t]{0.24\linewidth}
          \centering
          TIPS
        \end{minipage}%
        \hfill
        \begin{minipage}[t]{0.24\linewidth}
          \centering
          SigLIP2
        \end{minipage}%
        \hfill
        \begin{minipage}[t]{0.24\linewidth}
          \centering
          \tipsvtwo (ours)
        \end{minipage}%
    }%
    \vspace{-1mm}
	\caption{\small\textbf{PCA maps.} Comparing the first $3$ PCA components, using ViT-g models. 
    Images are forwarded at 1372 resolution for patch size 14 models (TIPS, TIPSv2) and at 1568 resolution for patch size 16 models (SigLIP2).
    \tipsvtwo produces smoother feature maps compared to other vision-language pretraining methods, with well-delineated objects.}
	\label{fig:qualitative_dense}
	\vspace{-3mm}
\end{figure}

\noindent\textbf{Image-only tasks.} In \cref{tab:sota_image_only}, \tipsvtwo achieves the best or second-best results in 7 out of 9 evaluations, and also achieves strong performance on the remaining 2 cases.
\tipsvtwo stands out as particularly strong in dense understanding tasks, improving by $+1.5$ in segmentation on PASCAL and by $-0.019$ on depth in NYUv2 over prior best.
However, our models do not perform as well in ImageNet classification, a result of our focus on optimizing for general purpose capability across a wide range of tasks.

\noindent\textbf{DINOv3 comparison.}
Additionally, we compare \tipsvtwo with the very recent DINOv3~\citep{siméoni2025dinov3}, presenting results in \cref{tab:dinov3}.
In order to make the zero-shot segmentation evals comparable, for this table we reproduce their setup, with the more expensive sliding window protocol from TCL~\citep{cha2023learning}.
Note that the DINOv3 teacher model is trained with 6$\times$ more parameters and 15$\times$ more images than the \tipsvtwo teacher, so we believe the fairest possible comparison is to use the largest model size in common between the two model families, which is ViT-L.
We report numbers on the subset of all released evals with matching protocols.
Notably, \tipsvtwo outperforms DINOv3 in 4 out of 6 evals, despite the data and model size disadvantage.


